\documentclass[11pt]{article}
\usepackage[margin=1in]{geometry}
\usepackage{amsmath}
\usepackage{booktabs}
\usepackage{tcolorbox}

\title{Tcolorbox Callout Stress Case}
\author{Test Corpus}
\date{}

\begin{document}
\maketitle

\begin{abstract}
This fixture tests a framed callout created with \texttt{tcolorbox}. The callout
contains ordinary prose and an inline equation so that conversion can be checked
for box geometry, text flow, and mathematical content.
\end{abstract}

\section{Decision Rule}
\label{sec:decision-rule}

Many academic appendices use boxed notes to identify assumptions or coding
decisions. The following callout records a compact rule for a synthetic sample:

\begin{tcolorbox}[
  title={Coding rule},
  colback=gray!6,
  colframe=black!65,
  boxrule=0.6pt,
  arc=1mm,
  left=6pt,
  right=6pt,
  top=6pt,
  bottom=6pt
]
A record enters the analytic sample when the eligibility score satisfies
$e_i \geq 3$ and the follow-up indicator is observed. Records below the threshold
remain in the tracking file for descriptive reporting.
\end{tcolorbox}

The paragraph after the callout should return to the normal text width. This
helps identify whether the box environment leaks spacing or indentation into
subsequent material.

\begin{table}[htbp]
\centering
\caption{Synthetic callout classification}
\label{tab:callout}
\begin{tabular}{lcc}
\toprule
Classification & Records & Included \\
\midrule
Eligible with follow-up & 138 & Yes \\
Eligible without follow-up & 21 & No \\
Below eligibility threshold & 46 & No \\
\bottomrule
\end{tabular}
\end{table}

Table~\ref{tab:callout} provides a standard float after the boxed material.

\end{document}
